\documentclass[11pt]{article}

\usepackage[margin=1.45in]{geometry}
\usepackage[T1]{fontenc}
\usepackage[utf8]{inputenc}
\usepackage{amsmath,amssymb}
\usepackage{enumitem}
\usepackage{booktabs}
\usepackage{array}
\usepackage[dvipsnames]{xcolor}
\usepackage[colorlinks=true,linkcolor=red,urlcolor=red,citecolor=red]{hyperref}

\setlist[itemize]{itemsep=2pt,topsep=3pt}
\setlist[enumerate]{itemsep=2pt,topsep=3pt}
\setlength{\emergencystretch}{3em}

\newcommand{\paper}[2]{\href{#2}{\textcolor{red}{#1}}}

\title{\textbf{Deep Learning in Finance}\\[4pt]
       \large Assignment 1 --- Research Paper Presentation and Implementation}
\author{Oualid Missaoui}
\date{}

\begin{document}

\maketitle

\section*{Overview}

Your first assignment is a \textbf{group presentation of a research paper},
delivered in class, together with a \textbf{GitHub repository that implements
the paper or a well-chosen segment of it}. This assignment is worth
\textbf{40\%} of your course grade.

The point of the exercise is not to summarize a paper. It is to get to the stage
where you could defend it: to know what the authors actually did, why the
architecture they chose was the right one for the problem, what the paper
establishes and what it merely asserts, and --- above all --- what happens when
you try to run it yourself. The implementation is where most of the learning
happens, and it is weighted accordingly.

\section*{Groups and Schedule}

The class splits into \textbf{five groups}. Presentations run from
\textbf{Week 3 through Week 7}, one group per session, at the start of class.
Each group is assigned a topic, which is lectured \emph{before} the
presentation, so you present on material the room has already seen.

Groups 1--5 fill the five sessions. If enrollment produces a \textbf{sixth or
seventh group}, those take the optional slots and double up with an existing
session --- two talks that week. Both optional topics are lectured well before
their slot, so those groups have the longest run-up of anyone.

\begin{center}
\begin{tabular}{@{}llp{5.6cm}c@{}}
\toprule
\textbf{Group} & \textbf{Presents} & \textbf{Topic} & \textbf{Lectured} \\
\midrule
Group 1 & Week 3 & Theoretical Foundations and the Virtue of Complexity & Week 1 \\
Group 2 & Week 4 & Transformers                                         & Week 2 \\
Group 3 & Week 5 & Meta-Learning and Regime Shifts                      & Week 3 \\
Group 4 & Week 6 & Deep Factor Models                                   & Week 4 \\
Group 5 & Week 7 & Generative Adversarial Networks                      & Week 4 \\
\midrule
\multicolumn{4}{@{}l}{\emph{Optional --- run only if enrollment produces a sixth or seventh group:}} \\
Group 6 & Week 6 & Diffusion Models                                     & Week 5 \\
Group 7 & Week 7 & Neural ODEs and Neural Point Processes               & Week 3 \\
\bottomrule
\end{tabular}
\end{center}

\noindent
Two topics are deliberately left uncovered. \textbf{Graph Neural Networks}
(Week 6) and \textbf{Bayesian Deep Networks for Factor Modeling} (Week 7) are
lectured too late in the term to leave a group a reasonable gap between the
lecture and its presentation slot. Both are lectured in the normal way; they
simply carry no presentation.

Group size depends on enrollment --- expect three to five students per group.
\textbf{Every member presents.} A group that sends one person to the front
while the rest sit down will be marked accordingly.

\section*{Choosing a Paper}

You have two routes:

\begin{enumerate}
  \item \textbf{Pick from the list} in the final section of this document. The
    PDFs are in \texttt{Assignments/papers/} in the course repository, with the
    exception noted there.
  \item \textbf{Propose your own.} Any paper is fair game provided it is a
    genuine research paper (peer-reviewed, or a serious preprint), it sits
    squarely within your group's topic, and \textbf{I approve it}. Email me the
    citation and a two-or-three sentence case for it. I am happy to approve
    papers that are not on my list --- a paper you are actually curious about
    will produce a better presentation than one you were handed. What I will
    turn down is a paper with no implementable core, a survey, or something so
    far from the topic that the rest of the class gains nothing.
\end{enumerate}

\noindent
\textbf{No two groups may present the same paper.} Choices are first-come,
first-served, so claim yours early.

\subsection*{Deadlines for selection}

\begin{itemize}
  \item \textbf{End of Week 1:} group rosters finalized and emailed to me.
  \item \textbf{End of Week 2:} paper choice submitted (or proposal approved).
    Group 1 should settle both as early in Week 1 as possible, since they
    present in Week 3.
\end{itemize}

\section*{Deliverable 1 --- The Presentation (10--15 minutes)}

Ten to fifteen minutes is short. It is not enough time to walk through a paper
section by section, and you should not try. Decide what the paper's one real
idea is and spend your time there. A rough shape that works:

\begin{itemize}
  \item \textbf{The problem (1--2 min).} What financial question is at stake,
    and why do the standard tools fail at it? Motivate it in the language of
    finance, not of machine learning.
  \item \textbf{The method (4--6 min).} The architecture and the objective.
    Write down the loss. Be precise about what is learned, what is fixed, and
    what the inputs and outputs actually are. If the paper's contribution is a
    modification of a standard architecture, make the modification visible.
  \item \textbf{The results (2--3 min).} The one or two tables or figures that
    carry the paper's claim. What is the baseline, and is it a fair one?
  \item \textbf{Your implementation (3--4 min).} What you built, what you
    reproduced, what you could not, and what you learned from the gap. This
    section is often the most interesting part of the talk --- do not let it
    get squeezed out.
  \item \textbf{Your critique (1--2 min).} Where is the paper vulnerable?
    Look-ahead bias, transaction costs, a suspiciously favorable sample
    period, a baseline chosen to lose, hyperparameters tuned on test data. Most
    papers have a soft spot. Finding it is a skill this course is trying to
    teach you.
\end{itemize}

\noindent
Expect questions, from me and from the class. Budget for them: aim your
prepared content at the lower end of the window. Slides in any format; send
them to me \textbf{24 hours before} your session.

\section*{Deliverable 2 --- The GitHub Repository}

Submit a link to a GitHub repository (public, or private with me added as a
collaborator) that implements the paper or a segment of it.

\subsection*{Scope: depth over breadth}

\textbf{You are not expected to reproduce a whole paper.} A faithful
implementation of one mechanism, evaluated honestly on a small universe, is
worth far more than a sprawling half-working attempt at everything. Legitimate
ways to cut scope:

\begin{itemize}
  \item One asset class, or one market, instead of the paper's full cross-section.
  \item A shorter sample period, or a smaller universe (e.g.\ 50--100 names).
  \item A reduced model --- fewer layers, smaller hidden dimension, fewer
    random features --- as long as the mechanism under study survives the cut.
  \item Simulated data where the paper's proprietary data is unobtainable. If
    the paper uses TAQ or CRSP and you cannot get it, build a simulator whose
    stylized facts you can state, and say clearly what that substitution costs
    you.
\end{itemize}

\noindent
State your scope decisions in the README. A clearly reasoned reduction is a
sign of judgment, not a shortfall. What is \emph{not} acceptable is cloning the
authors' repository and running it --- if an official implementation exists,
you may read it and should cite it, but the code you submit must be code you
wrote and understand.

\subsection*{What the repository must contain}

\begin{itemize}
  \item \textbf{A README} that states: the full citation of the paper; what you
    implemented and what you deliberately left out; how to install and run it;
    and what you found. Someone landing on the repo cold should understand the
    project in two minutes.
  \item \textbf{Reproducibility.} A \texttt{requirements.txt} or environment
    file, fixed random seeds, and a single documented entry point --- a script
    or a driver notebook --- that regenerates your headline result. I will try
    to run it.
  \item \textbf{Structure.} Reusable logic in modules; notebooks for narrative
    and figures. Not one 2{,}000-line notebook.
  \item \textbf{At least one figure or table} that either reproduces a result
    from the paper or is the honest analogue of one under your reduced scope,
    placed side by side with the paper's number where a comparison is possible.
  \item \textbf{Data handling.} Commit a small sample or a fetch script. Do not
    commit large vendor extracts.
  \item \textbf{A ``what didn't work'' section.} Required, not optional. The
    thing you could not get to train, the result that came out backwards, the
    number you cannot reconcile with the paper. Write it down. A repository
    reporting a failed reproduction with a clear account of the failure scores
    well; one quietly reporting success it did not achieve does not.
  \item \textbf{An AI usage note} (see below).
\end{itemize}

\noindent
\textbf{Due:} repository link emailed 24 hours before your session, with a
tagged commit (e.g.\ \texttt{v1.0}) marking what you are submitting. You are
welcome to keep working on it afterwards --- I grade the tag.

\section*{On Coding Agents and LLMs}

\textbf{Use them. Heavily. This is encouraged, not merely tolerated.}

Coding agents --- Claude Code, Codex, Cursor, and whatever has shipped by the
time you read this --- are now part of how quantitative research is actually
done, and you are expected to be fluent with them. Use an agent to scaffold
your repository, to port an equation into a working module, to write the
backtest harness, to track down the bug that has eaten your afternoon. There is
no credit here for doing by hand what a tool does better.

Likewise, \textbf{use LLMs to demystify the hard parts.} These are difficult
papers, often assuming machinery you have not met. When a paper asserts an
identity and moves on, ask. Ask for the derivation in smaller steps; ask for
the toy two-dimensional version of the construction; ask why the authors chose
this estimator over the obvious alternative; ask it to re-explain the same
object three different ways until one lands. Used this way, an LLM is the
patient co-reader that makes a hard paper tractable. I would rather you
attempted an ambitious paper with heavy LLM support than a safe one unaided.

Three conditions, and they are firm:

\begin{enumerate}
  \item \textbf{You own what you submit.} Anything in your repository or your
    slides, you can explain and defend on the spot. I will ask about specific
    lines. ``The agent wrote that'' is not an answer, and it is the one failure
    mode that this assignment is designed to expose. The implementation
    deliverable exists precisely because understanding survives it and
    cut-and-paste does not.
  \item \textbf{Verify everything.} Agents and LLMs produce confident,
    plausible, wrong output --- fabricated citations, subtly incorrect
    derivations, backtests with look-ahead baked in. Check the math against the
    paper. Check that your train/test split is what you think it is. A
    hallucinated result presented as the paper's is a serious error, and
    ``the model told me'' does not reduce your responsibility for it.
  \item \textbf{Disclose it.} Include a short \texttt{AI\_USAGE.md} (or a
    README section) describing which tools you used and for what --- a
    paragraph is plenty. This is not a penalty; nothing in your grade turns on
    how much you used them. It is professional practice, and it mirrors the
    disclosure in this repository's own README.
  \item \textbf{Do not let an agent write your talk.} Slides generated
    wholesale read exactly as such. The synthesis --- deciding what the
    paper's one real idea is, and what to cut --- is the part of this
    assignment that is yours.
\end{enumerate}

\section*{Grading}

\begin{center}
\begin{tabular}{@{}p{7.4cm}c@{}}
\toprule
\textbf{Component} & \textbf{Weight} \\
\midrule
Implementation --- correctness, reproducibility, and honesty of the repository & 40\% \\
Technical command of the paper --- the method explained precisely and correctly & 25\% \\
Presentation --- structure, clarity, use of the time, and handling of questions & 20\% \\
Critical assessment --- identifying the paper's real weaknesses & 15\% \\
\bottomrule
\end{tabular}
\end{center}

\noindent
As with the final project, \textbf{reproducing the paper's headline number is
not required for a strong grade.} Rigor, a correct account of the method, and a
clear-eyed explanation of why your results do or do not match carry more weight
than a matching Sharpe ratio. Groups are graded jointly, with an individual
adjustment where one member's contribution is plainly out of line with the rest.

\section*{Paper List}

\subsection*{Group 1 --- Theoretical Foundations and the Virtue of Complexity}

\begin{itemize}
  \item Bryan Kelly, Semyon Malamud, Kangying Zhou (2022) \paper{The Virtue of
    Complexity Everywhere}{https://papers.ssrn.com/sol3/papers.cfm?abstract_id=4166368},
    SSRN Working Paper 4166368. \emph{Not bundled --- download requires a free
    SSRN account.}
\end{itemize}

\subsection*{Group 2 --- Transformers}

\begin{itemize}
  \item Elliot L. Epstein, Rose Wang, Jaewon Choi, Markus
    Pelger (2025) \paper{Attention Factors for Statistical
    Arbitrage}{https://arxiv.org/pdf/2510.11616}.
\end{itemize}

\subsection*{Group 3 --- Meta-Learning and Regime Shifts}

\begin{itemize}
  \item Yicheng Wang, Sandro Claudio Lera (2025) \paper{Meta-Learning for
    Return Prediction in Shifting Market
    Regimes}{https://doi.org/10.1016/j.finmar.2025.101042}, Journal of
    Financial Markets, 101042.
\end{itemize}

\subsection*{Group 4 --- Deep Factor Models}

\begin{itemize}
  \item Yitong Duan, Lei Wang, Qizhong Zhang, Jian Li (2022)
    \paper{FactorVAE: A Probabilistic Dynamic Factor Model Based on Variational
    Autoencoder for Predicting Cross-Sectional Stock
    Returns}{https://doi.org/10.1609/aaai.v36i4.20369}, AAAI 36(4), 4468--4476.
\end{itemize}

\subsection*{Group 5 --- Generative Adversarial Networks}

\begin{itemize}
  \item Adil Rengim Cetingoz, Charles-Albert Lehalle (2025)
    \paper{Synthetic Data for Portfolios: A Throw of the Dice Will Never Abolish
    Chance}{https://arxiv.org/pdf/2501.03993}.
\end{itemize}

\subsection*{Group 6 (optional) --- Diffusion Models}

\begin{itemize}
  \item Zhuohan Wang, Carmine Ventre (2024) \paper{A
    Financial Time Series Denoiser Based on Diffusion
    Models}{https://doi.org/10.1145/3677052.3698649}, ICAIF.
\end{itemize}

\subsection*{Group 7 (optional) --- Neural ODEs and Neural Point Processes}

\begin{itemize}
  \item Jingsui Li, Wei Zhu, Zhang Chen, Chao Pei (2023)
    \paper{Financial Time Series Prediction via Neural Ordinary Differential
    Equations Approach}{https://doi.org/10.1109/CSIS-IAC60628.2023.10364112},
    CSIS-IAC.
\end{itemize}

\vspace{1em}
\noindent
Questions, paper proposals, and scope negotiations --- come early rather than
late, particularly about scope:

\begin{center}
\texttt{oualid.missaoui@baruch.cuny.edu} \quad or during office hours.
\end{center}

\end{document}
